% Métricas a reportar, esquema de validación (CV/holdout), umbrales de éxito y benchmark de negocio.

\section{Plan de modelado y validación}

\subsection{Métricas}
Para evaluar el rendimiento de los modelos se utilizarán métricas empleadas en \textcite[sección 4]{leborgne2022fraud}:
\begin{itemize}
\item \textit{ROC-AUC}: evalúa la capacidad del modelo para discriminar entre transacciones fraudulentas y legítimas.
\item \textit{Average Precision (AP)}: resume la relación entre precisión y \textit{recall}, y se calcula como
\(
    AP=\sum_n(R_n-R_{n-1})P_n,
\)
donde \(P_n\) y \(R_n\) representan la precisión y el \textit{recall} en cada punto.

\item \textit{Card Precision at \(k\)}: proporción de clientes fraudulentos entre los \(k\) clientes con mayor riesgo asignado por el modelo, donde \(k=100\) al igual que en el libro de \textcite{leborgne2022fraud}.
\end{itemize}

\subsection{Esquema de validación}

Debido al carácter temporal de las transacciones, la validación se realizará respetando el orden cronológico de los datos. Los modelos serán entrenados con transacciones anteriores y evaluados sobre períodos posteriores, evitando una partición aleatoria que pudiera introducir información futura en el entrenamiento.

Se considerará un retraso de siete días en la disponibilidad de las etiquetas de fraude, de manera consistente con las variables construidas durante el \textit{feature engineering}.

\subsection{Modelos y baseline}

Se establecerá como \textit{baseline} la regresión logística, a partir de la cual se evaluará si modelos de mayor complejidad producen mejoras en las métricas seleccionadas. 

\begin{nota}
    Se considerará una mejora respecto del benchmark cuando un modelo supere al baseline en la medida \textit{Card Precision at \(k\)}.
\end{nota}